% =============================================================================
% =============================================================================
% CAPÍTULO 1 --- INTRODUÇÃO GERAL
% =============================================================================
% =============================================================================
\chapter{Introdução geral}
\label{cap:introducao}
\todo{Ainda pendente de escrita}
% % % % % % % % % % % % % % % % % % % % % % % % % % % % % % % % % % % % % % %
% % BLOQUEADO - SEÇÃO CONTEXTUALIZAÇÃO, BIORREMEDIAÇÃO E PROBLEMA COMENTADA
% % % % % % % % % % % % % % % % % % % % % % % % % % % % % % % % % % % % % % %
% 
% % % % % % % % % % % % % % % % % % % % % % % % % % % % % % % % % % % % % % %
% % \section{Contextualização}
% % \label{sec:contextualizacao}
% % % % % % % % % % % % % % % % % % % % % % % % % % % % % % % % % % % % % % %
% 
% % \reaproveitar{Trazer do draft anterior os parágrafos sobre poluição ambiental, limitações de abordagens convencionais e avanço da era genômica (Seção 1.1 original).}
% 
% % \escrever{Parágrafo 1 --- Escala global da contaminação química:}
% % \begin{itemize}
% %     \item Poluentes prioritários (\glspl{pop}, hidrocarbonetos, pesticidas, metais pesados, contaminantes emergentes como PFAS e microplásticos).
% %     \item Impactos ecológicos e sanitários documentados.
% %     \item Citações de referência: relatórios da OMS, UNEP, EPA; artigos de revisão recentes sobre carga global de poluição.
% %     \todo{Inserir estatística atualizada sobre carga global de contaminação e citação}
% % \end{itemize}
% 
% % \escrever{Parágrafo 2 --- Limitações das abordagens físico-químicas convencionais:}
% % \begin{itemize}
% %     \item Custo, geração de resíduos secundários, insustentabilidade a longo prazo.
% %     \item Biorremediação como alternativa: uso de microrganismos, plantas e produtos metabólicos.
% %     \todo{Citação para justificar biorremediação como estratégia}
% % \end{itemize}
% 
% % \escrever{Parágrafo 3 --- Emergência das abordagens ômicas:}
% % \begin{itemize}
% %     \item \gls{ngs}, metagenômica, \glspl{mag}.
% %     \item Anotação funcional via identificadores ortólogos (\gls{ko}) permite inferir capacidades metabólicas sem cultivo.
% %     \item Convergência das ômicas (genômica, transcriptômica, metabolômica, proteômica) na caracterização do potencial biodegradativo.
% %     \todo{Citações sobre metagenômica ambiental e uso de KO em análises funcionais}
% % \end{itemize}
% 
% % \escrever{Parágrafo 4 --- Frameworks regulatórios como resposta institucional:}
% % \begin{itemize}
% %     \item Listagem de agências: \gls{iarc}, \gls{epa}, \gls{atsdr}, \gls{wfd}, \gls{psl}, \gls{epc}, \gls{conama}.
% %     \item Papel das listas prioritárias na definição de alvos de investigação e mitigação.
% % \end{itemize}
% 
% % % % % % % % % % % % % % % % % % % % % % % % % % % % % % % % % % % % % % %
% % % \section{Biorremediação como resposta multi-domínio}
% % % \label{sec:biorremediacao-multidominio}
% % % % % % % % % % % % % % % % % % % % % % % % % % % % % % % % % % % % % % %
% 
% % \nota{Esta seção é fundamental para a defesa do argumento de audiência ampla do trabalho. A biorremediação NÃO é nicho ambiental estrito; é família de aplicações que atravessa múltiplos setores. Redação sóbria e sem exagero, mas explicitando os domínios.}
% 
% % \escrever{Parágrafo introdutório:} Definir biorremediação em sentido funcional (transformação biológica de compostos químicos por sistemas vivos), não restrita a descontaminação ambiental.
% 
% % \escrever{Parágrafo(s) sobre domínios de aplicação --- cobrir de forma compacta:}
% % \begin{itemize}
% %     \item \textbf{Indústria de petróleo e gás}: recuperação avançada de óleo por microrganismos, biotratamento de águas produzidas e borras oleosas, biocorrosão.
% %     \item \textbf{Biocatálise farmacêutica e metabolismo de fármacos}: famílias enzimáticas (P450, monooxigenases, dehidrogenases, hidrolases) que mediam biotransformação tanto de poluentes quanto de fármacos; tratamento de efluentes farmacêuticos.
% %     \item \textbf{Degradação enzimática de plásticos}: PETase, MHETase; interseção com economia circular.
% %     \item \textbf{Fitorremediação e agronomia}: rizoremediação, recuperação de solos agrícolas.
% %     \item \textbf{Saúde pública, epidemiologia ambiental e regulação}: avaliação de risco ocupacional, cosmetovigilância, vigilância de contaminantes emergentes em água potável.
% % \end{itemize}
% 
% % \escrever{Parágrafo de fechamento:} A infraestrutura científica compartilhada entre estas comunidades (genes de degradação, rotas metabólicas, predição toxicológica, classificação regulatória) justifica uma plataforma unificada de análise.
% 
% % \todo{Citar 2-3 revisões recentes por domínio para embasar cada aplicação}
% 
% % % % % % % % % % % % % % % % % % % % % % % % % % % % % % % % % % % % % % %
% % % \section{Problema de pesquisa}
% % % \label{sec:problema}
% % % % % % % % % % % % % % % % % % % % % % % % % % % % % % % % % % % % % % %
% 
% % \reaproveitar{Trazer do draft anterior a formulação do gap (Seção 1.2 original), mas ajustar para refletir que o gap é sobre unificação regulatório-funcional, não apenas sobre integração de bases.}
% 
% % \escrever{Parágrafo 1 --- Formulação do gap:}
% % Não existe atualmente plataforma computacional que unifique, em arquitetura \textit{compound-centric}, anotação funcional (KEGG, HADEG), predição toxicológica (ToxCSM) e prioridade regulatória multi-jurisdicional, oferecendo interfaces complementares para diferentes perfis técnicos de usuário e cobertura efetiva de análises multi-ômicas.
% 
% % \escrever{Parágrafo 2 --- Consequências práticas do gap:}
% % \begin{itemize}
% %     \item Reconciliação manual de identificadores entre fontes heterogêneas.
% %     \item Desconexão entre análise funcional e relevância regulatória.
% %     \item Ausência de suporte à decisão para priorização baseada em risco.
% %     \item Dificuldade de reprodutibilidade quando a evidência é reconstruída caso a caso.
% % \end{itemize}
% 
% % \escrever{Parágrafo 3 --- Pergunta de pesquisa (frase única):}
% % Como projetar e implementar um ecossistema de bioinformática que integre evidência funcional, toxicológica e regulatória em torno de compostos prioritários, atendendo simultaneamente pesquisadores que operam por interfaces web e por pipelines programáticos?
% % % % % % % % % % % % % % % % % % % % % % % % % % % % % % % % % % % % % % %

% -----------------------------------------------------------------------------
\section{Objetivos}
\label{sec:objetivos}
% -----------------------------------------------------------------------------

\subsection{Objetivo geral}

Desenvolver o \biorempp{} --- \textit{Bioremediation Potential Profiling} ---,
um ecossistema integrado de bioinformática destinado à análise funcional de
dados ômicos e à integração de evidências biológicas, químicas, toxicológicas
e institucionais relacionadas a compostos prioritários para biorremediação. O
ecossistema compreende quatro componentes complementares e interoperáveis: o
banco de dados integrado BioRemPP, acompanhado por um pipeline reprodutível de
geração e validação; o BioRemPP Database Explorer
(\glsentryshort{dbx}), interface web destinada à exploração do banco; o
BioRemPP Web Service (\glsentryshort{bws}), destinado à análise de perfis
funcionais submetidos pelo usuário; e a BioRemPP Command-Line Interface
(\glsentryshort{cli}), destinada à execução programática e à integração com
pipelines automatizados.

\subsection{Objetivos específicos}


\begin{enumerate}

    \item Construir o banco de dados curado BioRemPP, cuja tabela central é gerada
    e validada por um pipeline reprodutível implementado em Snakemake e reúne
    compostos associados a nove códigos de referências institucionais. Integrar,
    por meio de identificadores interoperáveis, às relações funcionais
    e químicas do \gls{kegg}, incluindo um recorte de 20 vias de metabolismo e
    degradação de xenobióticos, às evidências especializadas de degradação do
    \gls{hadeg} e às predições toxicológicas do \gls{toxcsm}.

    \item Desenvolver o \gls{dbx}, uma interface web de acesso aberto e somente
    leitura para consulta e exploração do banco integrado, organizada em torno
    de compostos, classes químicas, genes, vias e informações toxicológicas, e
    complementada por análises guiadas.

    \item Implementar o \gls{bws}, um serviço web destinado à submissão de
    conjuntos de \glspl{ko} organizados por amostra, com módulos analíticos que
    projetem esses perfis sobre as relações do banco integrado e apoiem a
    caracterização exploratória, a comparação funcional e a formulação de
    hipóteses para estudos posteriores.

    \item Implementar o \gls{cli} e distribuí-lo por meio do \gls{pypi}, de
    modo a disponibilizar as operações do ecossistema para execução
    programática e incorporação a pipelines bioinformáticos automatizados, com
    modos operacionais destinados a entradas genômicas, transcriptômicas,
    metabolômicas e enzimáticas.

    \item Incorporar mecanismos alinhados aos princípios \gls{fair} por meio do
    uso consistente de identificadores interoperáveis, incluindo \gls{ko},
    \gls{cpd}, KEGG Reaction, \gls{ec}, \gls{cas}, \gls{chebi} e
    \gls{smiles}, do versionamento dos artefatos, do arquivamento em repositório
    com \gls{doi} persistente e da adoção de licenças abertas para o código
    (Apache~2.0) e para os dados (CC~BY~4.0).

    \item Examinar retrospectivamente a aplicabilidade do ecossistema por meio
    de três estudos publicados e revisados por pares que utilizaram dados ou
    relações do BioRemPP em contextos de caracterização funcional de novas
    espécies e subespécies e de integração de dados transcriptômicos e
    metatranscriptômicos.

\end{enumerate}

\newpage
% -----------------------------------------------------------------------------
\section{Organização da tese}
\label{sec:organizacao}
% -----------------------------------------------------------------------------

\todo{\textbf{Comentário:} Até a versão atual do manuscrito,
encontram-se redigidos os capítulos referentes ao banco de dados BioRemPP, ao
website de exploração do banco e ao website de análise dos dados submetidos
pelo usuário.}

A tese está organizada em sete capítulos. 
Após a introdução e a fundamentação teórica, cada componente
computacional é apresentado em um capítulo próprio, no qual são reunidos sua
motivação, seu escopo, sua metodologia e arquitetura, seus procedimentos de validação, seus
resultados e  discussão específica. O último capítulo retoma essas contribuições de maneira
integrada e discute o BioRemPP como um ecossistema computacional.

O primeiro capítulo apresenta o contexto científico e computacional que
motivou o desenvolvimento do BioRemPP. Nele são delimitados o problema de
pesquisa, a hipótese, os objetivos geral e específicos, as contribuições da
tese e a organização adotada para o manuscrito.

O capítulo dedicado à fundamentação teórica reúne a fundamentação teórica necessária à
interpretação dos componentes desenvolvidos. O capítulo aborda os fundamentos
da biorremediação, a contribuição das abordagens ômicas para a caracterização
do potencial funcional, o uso de grupos de ortologia e
identificadores químicos na integração de dados, as bases especializadas
relacionadas à degradação de contaminantes e as ferramentas computacionais
disponíveis nesse domínio. Também apresenta os princípios de organização,
interoperabilidade e reutilização de dados que orientam a adoção dos princípios
\gls{fair} no ecossistema.

O Capítulo~\ref{cap:comp1} apresenta o banco de dados integrado BioRemPP e o
processo computacional responsável por sua construção. O capítulo descreve as
fontes incorporadas, os critérios de curadoria, as relações estabelecidas entre
KOs, compostos, números EC, reações, vias, símbolos gênicos, descritores de
atividade enzimática, referências institucionais e predições toxicológicas.
Também detalha o workflow implementado em Snakemake, a organização do esquema
de dados, os procedimentos de verificação e validação e os resultados
correspondentes à versão 1.1.0 do banco.

O Capítulo~\ref{cap:comp2} descreve o BioRemPP Database Explorer
(\gls{dbx}), componente destinado à consulta e à exploração das relações
armazenadas no banco. O capítulo apresenta a arquitetura da aplicação, os
módulos de navegação, as análises guiadas e as
consultas SQL versionadas que acompanham as principais operações. Os resultados
demonstram como o componente organiza diferentes recortes do banco sem exigir
a submissão de dados externos ou a elaboração de consultas pelo usuário.

O Capítulo~\ref{cap:comp3} apresenta o BioRemPP Web Service
(\gls{bws}), componente destinado à análise interativa de conjuntos de KOs
organizados por amostra. O capítulo descreve o contrato de entrada, o fluxo de
validação e integração, a arquitetura em camadas e a definição declarativa dos
casos de uso. Também organiza os 56 casos de uso distribuídos em oito módulos,
que abrangem análises descritivas, multivariadas, relacionais, funcionais,
toxicológicas e de complementaridade entre amostras. Os resultados
demonstrativos utilizam perfis funcionais de amostras exemplo associadas na
literatura ao potencial de degradação de contaminantes e delimitam o alcance
interpretativo das análises baseadas na presença de KOs.

O capítulo dedicado à interface de linha de comando descreve a interface de linha de comando
(\gls{cli}) e sua incorporação a fluxos computacionais automatizados. O
capítulo apresenta a organização dos comandos, os contratos de entrada e
saída, o acesso programático às operações do ecossistema, os mecanismos de
configuração e os recursos destinados à execução reprodutível. Também discute
a relação entre a interface de linha de comando, o banco de dados, o
\gls{dbx} e o \gls{bws}, destacando os contextos nos quais a automação e a
integração com pipelines externos são necessárias.

O capítulo de conclusão integra as contribuições dos quatro componentes
e discute o BioRemPP como uma unidade computacional orientada à análise
funcional da biorremediação. O capítulo retoma os resultados centrais, examina
a complementaridade entre as formas de acesso ao banco e relaciona o
ecossistema a aplicações retrospectivas baseadas em estudos já publicados.
Também sistematiza as limitações associadas à cobertura das fontes, à qualidade
das anotações produzidas, ao uso de dados de presença e ausência e
à necessidade de validação experimental. Por fim, apresenta as perspectivas de
expansão do banco, incorporação de novas fontes, avaliação com usuários e
integração com outros fluxos de análise ômica.
